\section{How estimated components can change a sequential value}
\label{exec14:l5-remainder}

The exact score moments above assume the true conditional quantities. A
prototype must also make visible what happens when its fitted components are
wrong. This is particularly relevant when an assistant chooses a regression,
uses an estimated assignment model, and presents the resulting augmented
calculation as automatically protected against bias. Double robustness is a
property of a specified construction under specified conditions. It is not a
guarantee that any two imperfect components cancel each other.

Consider the same fixed risk policy, without retention loss or baseline
variation. Let $p_0$ be the observational probability of its first action and
$p_1(l)$ the probability of its required second action at state $l$. These are
probabilities of following the rule: when the rule requests zero, $p_1(l)$ is
the probability of action zero. Let $\widetilde p_0$ and
$\widetilde p_1(l)$ be fixed positive approximations, and define
\[
 \rho_0=\frac{p_0}{\widetilde p_0},\qquad
 \rho_1(l)=\frac{p_1(l)}{\widetilde p_1(l)}.
\]
Write $\widetilde m_g(l)=m_g(l)+d(l)$ and
$\widetilde q=V(g)+\eta$. The last quantity is an approximation to the
first-stage regression; it need not equal the average of
$\widetilde m_g$ under the true transition law. Replace the quantities in
equation~\eqref{eq:exec14-l5-score} by these fixed approximations and add
$V(g)$ to obtain an uncentered augmented value $\widetilde Z$. Conditioning
successively on the action history gives
\begin{align*}
 E[\widetilde Z]
 ={}&-\rho_0E_g[\rho_1(L_1)d(L_1)]\\
 &+\rho_0E_g[m_g(L_1)+d(L_1)-\widetilde q]
   +\widetilde q.
\end{align*}
Here $E_g$ averages over the state law after the regime's first action.
Subtracting $V(g)$ yields the exact identity
\begin{equation}
 E[\widetilde Z]-V(g)
 =(1-\rho_0)\eta
   +\rho_0E_g[(1-\rho_1(L_1))d(L_1)].
 \label{eq:exec14-l5-remainder}
\end{equation}
This is a project derivation for this restricted two-stage law, independently
checked by enumerating the observational histories
\citep{exec14_workflow_remainder}. It does not claim the general sequential
robustness theorem discussed by Luedtke and colleagues
\citep{sequentialdr2018}.

The identity separates two errors with different histories. Correct action
probabilities at both stages make the right side zero. Correct outcome
quantities at both stages also make it zero. Other combinations can make it
zero, but only by satisfying the displayed stage conditions. If $d=0$ while
both $\rho_0\neq1$ and $\eta\neq0$, correctness of the final outcome model
alone does not remove the first-stage product. Conversely, a correct
first-action probability does not remove a second-stage treatment/outcome
product. Fitting two separate models and calling one of them a critic has no
bearing on these conditions.

For a concrete failure, use the servicing law's $p_0=1/2$, approximate it by
$\widetilde p_0=2/5$, and set $\widetilde q=V(g)+3$. In the low-distress
state the rule requires action zero with true and approximated probability
$1/2$. In the distressed state it requires action one with true probability
$3/4$ and approximated probability $1/2$. Let $d(0)=0$ and $d(1)=5$.
Then $\rho_0=5/4$, and the second product averaged under the regime is
$(2/5)(1-3/2)5=-1$. The total bias is
\[
 (1-5/4)3+(5/4)(-1)=-2.
\]
The augmented mean becomes 92.8 instead of 94.8. Enumeration of the weighted
residual, intermediate regression and final scalar separately gives the same
answer. Restoring both assignment probabilities restores 94.8 despite these
outcome errors; restoring both outcome components also restores it despite
the assignment errors. These are population identities with fixed
approximations, not fitted-estimator performance results
\citep{exec14_workflow_remainder}.

There is also a quantitative diagnostic. With square-integrable components
under the regime's state law, Cauchy--Schwarz gives
\[
 |E[\widetilde Z]-V(g)|
 \leq |1-\rho_0|\,|\eta|
 +\rho_0\|1-\rho_1\|_{2,g}\|d\|_{2,g}.
\]
Lower bounds on approximated treatment probabilities are relevant because
they control the ratios; weak support can amplify small absolute propensity
errors. The bound does not supply observable values of the unknown errors.
It explains which approximations must be justified, and why predictive fit
alone need not settle the causal quantity.

For nuisance quantities trained on an independent sample, the fixed-component
calculation can be read conditionally on that training sample. A familiar
route to root-$n$ inference requires the bias products to be
$o_p(n^{-1/2})$, appropriate convergence of the centered score, and a valid
limit law for its sampling unit. Each requirement needs its own argument;
sample splitting alone proves none of them. Under the independent-account,
known-law construction above, the contrast's standard-error scale is
$\sqrt{667.52/n}$. A normal planning multiplier of 1.96 would require at
least 4007 accounts to make its planning half-width no greater than USD 0.8.
This is arithmetic using the known-law variance, not a claimed achieved
confidence interval, a selected trial size, or a bank precision commitment
\citep{exec14_workflow_remainder}. Dependence, estimated components and rare
histories can invalidate that planning scale. Increasing the account count
does not remove the fixed bias of two in the preceding example.

The assistant should therefore keep the causal target, fitted components,
training boundary, covariance unit and approximation status together. It can
flag a mistaken probability of action one when the policy requests zero,
propose checks of each stage, and compare a transparent fixed-regime
calculation. An unavailable rate argument remains an uncertainty limitation
even when the numerical value is finite. A report should explain that
limitation before turning the point estimate into a recommendation.

\section{Conditional panel comparisons must preserve cohort composition}
\label{exec14:l5-composition}

The preceding panel example used identical untreated trends across cohorts.
Conditional parallel trends allow a different situation: branches with the
same baseline covariates have comparable untreated changes, while adopting
and control branches have different covariate compositions. Pooling their
changes without standardization can then answer the wrong question even
though the conditional design is valid. Callaway and Sant'Anna's
identification representations average the conditional comparison over the
adopting group's covariate distribution \citep{callaway2021}.

Here is an independent stipulated two-period example in USD per account
\citep{exec14_workflow_remainder}. Let $X$ distinguish two branch markets,
and suppose three quarters of adopting branches have $X=1$, compared with
one quarter of never-adopting controls. At $X=0$, the untreated change is two
and the adoption effect is four. At $X=1$, those quantities are ten and
eight. Within each market the untreated change is identical for adopting
and control branches; no anticipation and stable composition are stipulated.
The target is the adopting cohort's average treatment effect,
\[
 \operatorname{ATT}=\tfrac14(4)+\tfrac34(8)=7.
\]
The adopting cohort's observed average change is
$\tfrac14(2+4)+\tfrac34(10+8)=15$. The pooled control change is
$\tfrac34(2)+\tfrac14(10)=4$. Their unadjusted difference is eleven,
overstating the target by four. The difference is caused by unequal market
composition, not by a violation of the stipulated conditional trends.

Outcome standardization instead computes the control change under the
adopting composition: $\tfrac14(2)+\tfrac34(10)=8$. Subtracting it from
fifteen restores seven. Weighting supplies a second representation. Suppose
the two cohorts have equal population mass, so the propensity to belong to
the adopting cohort is $e(0)=1/4$ and $e(1)=3/4$. Control weights proportional
to $e(X)/\{1-e(X)\}$ are $1/3$ and three. Their control-group average is
one, and their weighted change is
\[
 \tfrac34\left(\tfrac13\right)2
 +\tfrac14(3)10=8.
\]
The weights reconstruct the adopting covariate distribution. They are
cohort-comparison weights, rather than the probabilities of following a
two-action account policy. The corresponding target and observation clocks
must remain distinct when an assistant selects an API.

If market $X=1$ has no untreated controls, this representation cannot be
evaluated from the stipulated comparison without extrapolation or another
restriction. Dropping that market changes the target from seven to four;
silently describing the retained-market result as the original ATT would
replace the business question. Matching, regression and weighting should
therefore be compared on their supported target, not merely on whether each
returns a coefficient. Adding a covariate measured after platform adoption
can also change the scientific argument; the baseline market variable here
does not justify adjustment for any convenient post-adoption metric.

For an illustrative deployment cost of six per adopting account, the
identified gross effect of seven leaves a net margin of one. The pooled
difference of eleven would suggest five. These are planning values under
stipulated truth and cost, with no fitted uncertainty or bank ROI claim
\citep{exec14_workflow_remainder}. They show why composition and support checks
matter to the investment proposal: a target error can dominate a modest
economic margin even when every arithmetic step and generated explanation
appears convincing.
\label{exec14:l5-depth-end}
